\ifdefined\gkver\else\def\gkver{student}\fi
\documentclass[\gkver]{gaokaozhenti}
\begin{document}

\chapter{2018年海南理综}
%% paperType: 理综物理部分

\begin{enumerate}
\item
%% number: 1
%% paperName: 2018年海南理综第 1 题 4 分
%% typeId: 1
%% type: 单选题
%% chapter: 直线运动
%% point: 自由落体运动
%% method: 近似与估算
%% score: 4
%% degree: 850
%% duplicateId: 0
%% body:
一攀岩者以 $1\Ums$ 的速度匀速向上攀登，途中碰落了岩壁上的石块，石块自由下落。$3\Us$ 后攀岩者听到石块落地的声音，此时他离地面的高度约为\xzanswer{C}

\fourchoices[answer=C]
{$10\Um$}
{$30\Um$}
{$50\Um$}
{$70\Um$}

%% answer: C
%% memo:
\memoanswer{
	本题考查自由落体运动遵循的规律及其相关的知识点，同时考查了估算能力。

	根据自由落体运动遵循的规律可知，石块自由下落的高度约为 $h = \frac{1}{2}gt^{2} = 44.1\Um$；用时 $3\Us$，攀岩者爬了 $3\Um$，所以距离地面高度约为 $h \approx 44.1 + 3 = 47.1\Um$，则考虑到空气阻力和声速的影响，他离地面的高度约为 $50\Um$，选项 C 正确。
}

\item
%% number: 2
%% paperName: 2018年海南理综第 2 题 4 分
%% typeId: 1
%% type: 单选题
%% chapter: 曲线运动
%% point: 天体运动
%% method: 无
%% score: 4
%% degree: 750
%% duplicateId: 0
%% body:
土星与太阳的距离是火星与太阳距离的 6 倍多。由此信息可知\xzanswer{B}

\fourchoices[answer=B]
{土星的质量比火星的小}
{土星运行的速率比火星的小}
{土星运行的周期比火星的小}
{土星运行的角速度大小比火星的大}

%% answer: B
%% memo:
\memoanswer{
	本题考查圆周运动及万有引力定律相关的知识点。

	根据万有引力提供向心力 $G\frac{Mm}{r^{2}} = m\frac{v^{2}}{r}$，得 $v = \sqrt{\frac{GM}{r}}$，即 $r$ 越大速度越小，但是不比较土星和地球质量大小，故 A 错误，B 正确；由 $G\frac{Mm}{r^{2}} = m\frac{4\pi^{2}}{T^{2}}r$，得 $T = \sqrt{\frac{4\pi^{2}r^{3}}{GM}}$，即 $r$ 越大周期越大，故 C 错误；由 $G\frac{Mm}{r^{2}} = m\omega^{2}r$，得 $\omega = \sqrt{\frac{GM}{r^{3}}}$，即 $r$ 越大角速度越小，故 D 错误。
}

\item
%% number: 3
%% paperName: 2018年海南理综第 3 题 4 分
%% typeId: 1
%% type: 单选题
%% chapter: 磁场
%% point: 安培力
%% method: 无
%% score: 4
%% degree: 850
%% duplicateId: 0
%% body:
如图，一绝缘光滑固定斜面处于匀强磁场中，磁场的磁感应强度大小为 $B$，方向垂直于斜面向上，通有电流 $I$ 的金属细杆水平静止在斜面上。若电流变为 $0.5I$，磁感应强度大小变为 $3B$，电流和磁场的方向均不变，则金属细杆将\xzanswer{A}

\onepicture[w=4.4cm]{\includesvg{figs/03}}

\fourchoices[answer=A]
{沿斜面加速上滑}
{沿斜面加速下滑}
{沿斜面匀速上滑}
{仍静止在斜面上}

%% answer: A
%% memo:
\memoanswer{
	本题是一道关于安培力的共点力的平衡问题及牛顿第二定律相关知识。

	最初金属细杆受到三个力作用，且合力为零，如图所示：

	\onepicture[w=3.4cm]{figs/03_an.jpg}

	由平衡可知，安培力 $F = BIL = mg\sin\theta$，若电流变为 $0.5I$，磁感应强度大小变为 $3B$，则安培力 $F_{1} = 3B \times \frac{1}{2}IL = 1.5mg\sin\theta$，根据牛顿第二定律，$F_{1} - mg\sin\theta = ma$，故金属细杆以 $a = 0.5g\sin\theta$ 的加速度沿着斜面加速上滑，故 A 正确。
}

\item
%% number: 4
%% paperName: 2018年海南理综第 4 题 4 分
%% typeId: 1
%% type: 单选题
%% chapter: 原子和原子核
%% point: 半衰期
%% method: 无
%% score: 4
%% degree: 850
%% duplicateId: 0
%% body:
已知 \ce{^{234}_{90}Th} 的半衰期为 24 天。$4\Ug$ \ce{^{234}_{90}Th} 经过 72 天还剩下\xzanswer{B}

\fourchoices[answer=B]
{$0$}
{$0.5\Ug$}
{$1\Ug$}
{$1.5\Ug$}

%% answer: B
%% memo:
\memoanswer{
	本题考查核衰变的半衰期的规律及其相关的知识点。

	由衰变公式 $m' = m\left(\frac{1}{2}\right)^{\frac{t}{T}}$，知 $m' = 4\left(\frac{1}{2}\right)^{\frac{72}{24}} = 4\left(\frac{1}{2}\right)^{3} = 0.5\Ug$，故 B 正确。
}

\item
%% number: 5
%% paperName: 2018年海南理综第 5 题 4 分
%% typeId: 1
%% type: 单选题
%% chapter: 运动和力
%% point: 动量和能量
%% method: 无
%% score: 4
%% degree: 650
%% duplicateId: 0
%% body:
如图，用长为 $l$ 的轻绳悬挂一质量为 $M$ 的沙箱，沙箱静止。一质量为 $m$ 的弹丸以速度 $v$ 水平射入沙箱并留在其中，随后与沙箱共同摆动一小角度。不计空气阻力。对子弹射向沙箱到与其共同摆过一小角度的过程\xzanswer{B}

\onepicture[h=3.6cm]{\includesvg{figs/05}}

\fourchoices[answer=B]
{若保持 $m$、$v$、$l$ 不变，$M$ 变大，则系统损失的机械能变小}
{若保持 $M$、$v$、$l$ 不变，$m$ 变大，则系统损失的机械能变小}
{若保持 $M$、$m$、$l$ 不变，$v$ 变大，则系统损失的机械能变大}
{若保持 $M$、$m$、$v$ 不变，$l$ 变大，则系统损失的机械能变大}

%% answer: B
%% memo:
\memoanswer{
	本题考查机械能守恒定律及动量守恒定律的相关知识点。

	分析易得只有子弹射入沙箱并停留在沙箱中这个过程中系统的机械能才会损失，由动量守恒得，$mv = (m + M)v_{\text{共}}$，则系统损失的机械能

	$\Delta E = \frac{1}{2}mv^{2} - \frac{1}{2}(m + M)v_{\text{共}}^{2} = \frac{1}{2}mv^{2} - \frac{1}{2}(m + M)\left(\frac{mv}{m + M}\right)^{2} = \frac{mMv^{2}}{2(m + M)}$

	对于 A 选项，由 $\Delta E = \frac{mMv^{2}}{2(m + M)} = \frac{mv^{2}}{2\left(\frac{m}{M} + 1\right)}$ 可知若保持 $m$、$v$、$l$ 不变，$M$ 越大则系统损失的机械能越大，故 A 错误；

	对于 B 选项，由 $\Delta E = \frac{mMv^{2}}{2(m + M)} = \frac{Mv^{2}}{2\left(\frac{M}{m} + 1\right)}$ 可知若保持 $M$、$v$、$l$ 不变，$m$ 变大则系统损失的机械能变大，故 B 错误；

	对于 C 选项，由 $\Delta E = \frac{mMv^{2}}{2(m + M)}$ 可知若保持 $M$、$m$、$l$ 不变，$v$ 变大则系统损失的机械能变大，故 C 正确；

	对于 D 选项，子弹与沙箱共同向上摆动一小角度的过程系统机械能守恒，故 D 错误。

	（本题原卷及材料解析均为 C，而平台 JSON 答案标 B，二者不一致；答案命令按平台数据取 B，解析照原卷如实保留。）
}

\item
%% number: 6
%% paperName: 2018年海南理综第 6 题 4 分
%% typeId: 1
%% type: 单选题
%% chapter: 功和能
%% point: 功率
%% method: 无
%% score: 4
%% degree: 750
%% duplicateId: 0
%% body:
某大瀑布的平均水流量为 $5900\Umcs$，水的落差为 $50\Um$。已知水的密度为 $1.00\times10^{3}\Ukgmc$。在大瀑布水流下落过程中，重力做功的平均功率约为\xzanswer{D}

\fourchoices[answer=D]
{$3\times10^{6}\UW$}
{$3\times10^{7}\UW$}
{$3\times10^{8}\UW$}
{$3\times10^{9}\UW$}

%% answer: D
%% memo:
\memoanswer{
	本题考查平均功率及平均水流量的相关知识点，同时考查了估算能力。

	由平均功率定义得

	$P = \frac{W}{t} = \frac{mgh}{t} = \frac{Qt\rho gh}{t} = Q\rho gh \approx 5900 \times 1 \times 10^{3} \times 10 \times 50 \approx 6000 \times 1 \times 10^{3} \times 10 \times 50 = 3 \times 10^{9}\UW$

	故 D 正确。
}

\item
%% number: 7
%% paperName: 2018年海南理综第 7 题 5 分
%% typeId: 2
%% type: 多选题
%% chapter: 电磁感应
%% point: 切割磁感线电动势
%% method: 无
%% score: 5
%% degree: 650
%% duplicateId: 0
%% body:
如图，在磁感应强度大小为 $B$ 的匀强磁场中，有一面积为 $S$ 的矩形单匝闭合导线 abcd，ab 边与磁场方向垂直，线框的电阻为 $R$。使线框以恒定角速度 $\omega$ 绕过 ad、bc 中点的轴旋转。下列说法正确的是\xzanswer{AC}

\onepicture[w=4.0cm]{\includesvg{figs/07}}

\fourchoices[answer=AC]
{线框 abcd 中感应电动势的最大值是 $BS\omega$}
{线框 abcd 中感应电动势的有效值是 $BS\omega$}
{线框平面与磁场方向平行时，流经线框的电流最大}
{线框平面与磁场方向垂直时，流经线框的电流最大}

%% answer: AC
%% memo:
\memoanswer{
	本题考查了交变电流的最大值、有效值和瞬时值的相关知识点。

	一个单匝线圈在匀强磁场中旋转，当从中性面开始计时，产生的正弦式交变电流电动势的瞬时值表达式为：$e = E_{m}\sin\theta = E_{m}\sin\omega t$。

	故感应电动势的最大值 $E_{m} = BS\omega$，有效值 $E = \frac{E_{m}}{\sqrt{2}}$，故 A 正确，B 错误；

	当 $\theta = 90^\circ$ 时，即线框平面与磁场方向平行时，产生的感应电动势最大，流经线框的电流最大，故 C 正确；线框平面与磁场方向垂直时，产生的感应电动势为零，感应电流为零，故 D 错误。
}

\item
%% number: 8
%% paperName: 2018年海南理综第 8 题 5 分
%% typeId: 1
%% type: 单选题
%% chapter: 运动和力
%% point: 传送带
%% method: 无
%% score: 5
%% degree: 550
%% duplicateId: 0
%% body:
如图~\ref{2018海南理综08a}，一长木板静止于光滑水平桌面上，$t = 0$ 时，小物块以速度 $v_{0}$ 滑到长木板上，图~\ref{2018海南理综08b}为物块与木板运动的 $v-t$ 图像，图中 $t_{1}$、$v_{0}$、$v_{1}$ 已知。重力加速度大小为 $g$。由此可求得\xzanswer{BC}

\twopicture[labelA=2018海南理综08a, labelB=2018海南理综08b, wA=4.5cm, wB=6.0cm]
{figs/08a.png}{figs/08b.png}

\fourchoices[answer=BC]
{木板的长度}
{物块与木板的质量之比}
{物块与木板之间的动摩擦因数}
{从 $t = 0$ 开始到 $t_{1}$ 时刻，木板获得的动能}

%% answer: BC
%% memo:
\memoanswer{
	本题考查了 $v-t$ 图像与牛顿第二定律综合运用，滑块模型等。

	A．根据题意只能求出 A、B 的相对位移，不知道 B 最终停在哪里，无法求出木板的长度，故 A 不能够求解出；

	由图象的斜率表示加速度求出长木板的加速度为 $a_{A} = \frac{v_{1}}{t_{1}}$，小物块的加速度 $a_{B} = \frac{v_{0} - v_{1}}{t_{1}}$，

	根据牛顿第二定律得：$\mu mg = Ma_{A}$，$\mu mg = ma_{B}$，解得：$\frac{m}{M} = \frac{v_{1}}{v_{0} - v_{1}}$，$\mu = \frac{v_{0} - v_{1}}{gt_{1}}$，故 B 和 C 能够求解出；

	D．木板获得的动能 $E_{kA} = \frac{1}{2}Mv_{1}^{2} = \frac{1}{2}mv_{1}(v_{0} - v_{1})$，题目 $t_{1}$、$v_{0}$、$v_{1}$ 已知，但是 $M$、$m$ 不知道，故 D 不能够求解出。
}

\item
%% number: 9
%% paperName: 2018年海南理综第 9 题 5 分
%% typeId: 2
%% type: 多选题
%% chapter: 电场
%% point: 电场线
%% method: 无
%% score: 5
%% degree: 450
%% duplicateId: 0
%% body:
如图，a、b、c、d 为一边长为 $l$ 的正方形的顶点。电荷量均为 $q$（$q > 0$）的两个点电荷分别固定在 a、c 两点，静电力常量为 $k$。不计重力。下列说法正确的是\xzanswer{AD}

\onepicture[w=4.2cm]{\includesvg{figs/09}}

\fourchoices[answer=AD]
{b 点的电场强度大小为 $\frac{\sqrt{2}kq}{l^{2}}$}
{过 b、d 点的直线位于同一等势面上}
{在两点电荷产生的电场中，ac 中点的电势最低}
{在 b 点从静止释放的电子，到达 d 点时速度为零}

%% answer: AD
%% memo:
\memoanswer{
	本题考查电场强度以及等量同种电荷的电场、电势图的相关知识点。

	由图可知 b 点的电场 $E = E_{a}\cos 45^\circ + E_{b}\cos 45^\circ = \frac{\sqrt{2}kq}{l^{2}}$，故 A 正确。

	\onepicture[w=4.0cm]{figs/09_an1.jpg}

	沿着电场线电势逐渐降低，而等量正点电荷的电场与电势图如下，由图可知过 b、d 点的直线不在同一等势面上，故 B、C 错误；由对称性可知，b、d 点电势相同，故电子在 b、d 点电势能相同，即动能也相同，都为 0，故 D 正确。

	\onepicture[w=4.5cm]{figs/09_an2.jpg}
}

\item
%% number: 10
%% paperName: 2018年海南理综第 10 题 5 分
%% typeId: 2
%% type: 多选题
%% chapter: 电路
%% point: 动态分析
%% method: 无
%% score: 5
%% degree: 450
%% duplicateId: 0
%% body:
如图，三个电阻 $R_{1}$、$R_{2}$、$R_{3}$ 的阻值均为 $R$，电源的内阻 $r < R$，c 为滑动变阻器的中点。闭合开关后，将滑动变阻器的滑片由 c 点向 a 端滑动，下列说法正确的是\xzanswer{CD}

\onepicture[w=5.2cm]{figs/10.png}

\fourchoices[answer=CD]
{$R_{2}$ 消耗的功率变小}
{$R_{3}$ 消耗的功率变大}
{电源输出的功率变大}
{电源内阻消耗的功率变大}

%% answer: CD
%% memo:
\memoanswer{
	本题考查闭合电路欧姆定律的综合动态分析问题的相关知识点。

	把等效电路画出，如下：

	\onepicture[w=5.0cm]{figs/10_an1.jpg}

	设 $R_{ca} = R_{x}$，$R_{cd} = R_{\text{滑}} - R_{x}$，则 $R_{\text{外}} = R + \frac{(R + R_{x})(R + R_{\text{滑}} - R_{x})}{R + R + R_{\text{滑}}}$

	当 $R_{x} = \frac{R_{\text{滑}}}{2}$ 时，$R_{\text{外}}$ 有最大值；当滑动变阻器的滑片由 c 点向 a 端滑动时，$R_{ca}$ 减小，$R_{cd}$ 增大，易得 $R_{\text{外}}$ 减小，

	$R_{\text{总}} = R_{\text{外}} + r$ 减小，$I_{1} = \frac{E}{R_{\text{总}}}$ 增大，故电源内阻消耗的功率 $P_{r} = I_{1}^{2}r$ 增大，故 D 正确；

	$U_{R_{3} + R_{cb}} = U_{R_{2} + R_{ca}} = E - I_{1}(r + R_{1})$ 减小，

	$I_{3} = \frac{U_{R_{3} + R_{cb}}}{R_{3} + R_{cb}}$ 减小，故 $P_{R3} = I_{3}^{2}R_{3}$ 减小，故 B 错误；

	而 $I_{2} = I_{1} - I_{3}$ 增大，故 $P_{R2} = I_{2}^{2}R_{2}$ 增大，故 A 错误；

	根据电源输出功率与 $R_{\text{外}}$ 的关系图可知，当 $R_{\text{外}} > r$ 时，$R_{\text{外}}$ 减小则电源输出功率越大，故 C 正确。

	\onepicture[w=5.0cm]{figs/10_an2.jpg}
}

\item
%% number: 11
%% paperName: 2018年海南理综第 11 题 6 分
%% typeId: 4
%% type: 实验题
%% chapter: 直线运动
%% point: 自由落体运动
%% method: 无
%% score: 6
%% degree: 650
%% duplicateId: 0
%% body:
学生课外实验小组使用如图所示的实验装置测量重力加速度大小。实验时，他们先测量分液漏斗下端到水桶底部的距离 $s$；然后使漏斗中的水一滴一滴地下落，调整阀门使水滴落到桶底发出声音的同时，下一滴水刚好从漏斗的下端滴落；用秒表测量第 1 个水滴从漏斗的下端滴落至第 $n$ 个水滴落到桶底所用的时间 $t$。

\begin{enumerate}
	\item
	重力加速度大小可表示为 $g = $\tkanswer[2cm]{$\frac{2n^{2}s}{t^{2}}$}（用 $s$、$n$、$t$ 表示）；
	\item
	如果某次实验中，$s = 0.90\Um$，$n = 30$，$t = 13.0\Us$，则测得的重力加速度大小 $g = $\tkanswer[1.5cm]{9.6}$\Umsq$；（保留 2 位有效数字）
	\item
	写出一条能提高测量结果准确程度的建议：\tkanswer[4cm]{适当增大 $n$（或多次测量取平均值）}。
\end{enumerate}

\onepicture[h=6.5cm, align=right]{figs/11.png}

%% answer: 
\jdanswer{
\begin{enumerate}
	\item
	$\frac{2n^{2}s}{t^{2}}$

	\item
	9.6

	\item
	适当增大 $n$（或多次测量取平均值）

\end{enumerate}
}
%% memo:
\memoanswer{
	本题考查一个利用自由落体运动规律测量重力加速度的实验。

	（1）已知第 1 个水滴从漏斗的下端滴落至第 $n$ 个水滴落到桶底所用的时间 $t$，所以一个水滴从漏斗的下端滴落到桶底所用的时间 $\Delta t = \frac{t}{n}$，所以 $s = \frac{1}{2}g(\Delta t)^{2} = \frac{1}{2}g\left(\frac{t}{n}\right)^{2}$，故 $g = \frac{2n^{2}s}{t^{2}}$。

	（2）$g = \frac{2n^{2}s}{t^{2}} = \frac{2 \times 30^{2} \times 0.90}{13.0^{2}}\Umsq = 9.6\Umsq$。

	（3）提高精度可以提高测量时间的精度和距离的精度，也可以适当增大 $n$，或多测几次取平均。
}

\item
%% number: 12
%% paperName: 2018年海南理综第 12 题 12 分
%% typeId: 4
%% type: 实验题
%% chapter: 电路
%% point: 测电动势与内阻
%% method: 无
%% score: 12
%% degree: 450
%% duplicateId: 0
%% body:
某同学利用图~\ref{2018海南理综12a}中的电路测量电流表 A 的内阻 $R_{A}$（约为 $5\UO$）和直流电源的电动势 $E$（约为 $10\UV$）。图中 $R_{1}$ 和 $R_{2}$ 为电阻箱，S$_{1}$ 和 S$_{2}$ 为开关。已知电流表的量程为 $100\UmA$，直流电源的内阻为 $r$。

\begin{enumerate}
	\item
	断开 S$_{2}$，闭合 S$_{1}$，调节 $R_{1}$ 的阻值，使 A 满偏；保持 $R_{1}$ 的阻值不变，闭合 S$_{2}$，调节 $R_{2}$，当 $R_{2}$ 的阻值为 $4.8\UO$ 时 A 的示数为 $48.0\UmA$。忽略 S$_{2}$ 闭合后电路中总电阻的变化，经计算得 $R_{A} = $\tkanswer[1.5cm]{5.2}$\UO$；（保留 2 位有效数字）
	\item
	保持 S$_{1}$ 闭合，断开 S$_{2}$，多次改变 $R_{1}$ 的阻值，并记录电流表的相应示数。若某次 $R_{1}$ 的示数如图~\ref{2018海南理综12b}所示，则此次 $R_{1}$ 的阻值为\tkanswer[1.5cm]{148.2}$\UO$；
	\item
	利用记录的 $R_{1}$ 的阻值和相应的电流表示数 $I$，作出 $I^{-1}-R_{1}$ 图线，如图~\ref{2018海南理综12c}所示。用电池的电动势 $E$、内阻 $r$ 和电流表内阻 $R_{A}$ 表示 $I^{-1}$ 随 $R_{1}$ 变化的关系式为 $I^{-1} = $\tkanswer[2.5cm]{$\frac{R_{A} + r}{E} + \frac{1}{E}R_{1}$}。利用图~\ref{2018海南理综12c}可求得 $E = $\tkanswer[1.5cm]{9.1}$\UV$。（保留 2 位有效数字）
\end{enumerate}

\threepicture[labelA=2018海南理综12a, labelB=2018海南理综12b, labelC=2018海南理综12c, wA=5.0cm, wB=4.6cm, wC=5.2cm, align=right]
{figs/12a.png}{figs/12b.png}{figs/12c.png}

%% answer: 
\jdanswer{
\begin{enumerate}
	\item
	5.2

	\item
	148.2

	\item
	$\frac{R_{A} + r}{E} + \frac{1}{E}R_{1}$，9.1

\end{enumerate}
}
%% memo:
\memoanswer{
	本题考查一个半偏法测表头内阻的实验，附带测量电源的电动势。

	（1）因为忽略 S$_{2}$ 闭合后电路中总电阻的变化，总电流还是为 $100\UmA$，表头与 $R_{2}$ 并联，电流按电阻成反比分配，即：$R_{A}:R_{2} = I_{2}:I_{A} = 52:48$，易知：$R_{A} = 5.2\UO$。

	（2）由图可知 $R_{1} = 1 \times 100\UO + 4 \times 10\UO + 8 \times 1\UO + 2 \times 0.1\UO = 148.2\UO$。

	（3）根据闭合电路欧姆定律，$I = \frac{E}{R_{A} + r + R_{1}}$

	易得 $\frac{1}{I} = \frac{R_{1}}{E} + \frac{R_{A} + r}{E}$

	由斜率 $k = \frac{1}{E} = \frac{22.5 - 12}{196 - 100}$，可求出 $E = \frac{1}{k} = \frac{196 - 100}{22.5 - 12}\UV \approx 9.1\UV$。
}

\item
%% number: 13
%% paperName: 2018年海南理综第 13 题 10 分
%% typeId: 6
%% type: 计算题
%% chapter: 磁场
%% point: 带电粒子在磁场中的运动
%% method: 无
%% score: 10
%% degree: 650
%% duplicateId: 0
%% body:
如图，圆心为 O、半径为 $r$ 的圆形区域外存在匀强磁场，磁场方向垂直于纸面向外，磁感应强度大小为 $B$。P 是圆外一点，OP $= 3r$。一质量为 $m$、电荷量为 $q$（$q > 0$）的粒子从 P 点在纸面内垂直于 OP 射出。已知粒子运动轨迹经过圆心 O，不计重力。求：

\begin{enumerate}
	\item
	粒子在磁场中做圆周运动的半径；
	\item
	粒子第一次在圆形区域内运动所用的时间。
\end{enumerate}

\onepicture[w=6.5cm, align=right]{figs/13.png}

%% answer: 
\jdanswer{
\begin{enumerate}
	\item
	$R = \frac{4}{3}r$

	\item
	$t = \frac{3m}{2qB}$

\end{enumerate}
}
%% memo:
\memoanswer{
	本题考查在匀强磁场中的匀速圆周运动及其相关的知识点，意在考查考生灵活运用相关知识解决问题的能力。

	（1）找圆心，画轨迹，求半径。

	\onepicture[w=5.5cm]{figs/13_an.jpg}

	设粒子在磁场中运动半径为 $R$，由几何关系得：$R + \sqrt{R^{2} + r^{2}} = 3r$

	易得 $R = \frac{4}{3}r$。

	（2）设进入磁场时速度的大小为 $v$，由洛伦兹力公式和牛顿第二定律有

	$qvB = \frac{mv^{2}}{R}$

	进入圆形区域，带电粒子做匀速直线运动，则

	$2r = vt$

	联立解得 $t = \frac{3m}{2qB}$。
}

\item
%% number: 14
%% paperName: 2018年海南理综第 14 题 16 分
%% typeId: 6
%% type: 计算题
%% chapter: 运动和力
%% point: 动量和能量
%% method: 无
%% score: 16
%% degree: 350
%% duplicateId: 0
%% body:
如图，光滑轨道 PQO 的水平段 QO $= \frac{h}{2}$，轨道在 O 点与水平地面平滑连接。一质量为 $m$ 的小物块 A 从高 $h$ 处由静止开始沿轨道下滑，在 O 点与质量为 $4m$ 的静止小物块 B 发生碰撞。A、B 与地面间的动摩擦因数均为 $\mu = 0.5$，重力加速度大小为 $g$。假设 A、B 间的碰撞为完全弹性碰撞，碰撞时间极短。求：

\begin{enumerate}
	\item
	第一次碰撞后瞬间 A 和 B 速度的大小；
	\item
	A、B 均停止运动后，二者之间的距离。
\end{enumerate}

\onepicture[w=6.5cm, align=right]{figs/14.png}

%% answer: 
\jdanswer{
\begin{enumerate}
	\item
	$\frac{3}{5}\sqrt{2gh}$，$\frac{2}{5}\sqrt{2gh}$

	\item
	$\frac{26}{125}h$

\end{enumerate}
}
%% memo:
\memoanswer{
	本题主要考查机械能、匀变速直线运动规律、动量守恒定律、能量守恒定律及其相关的知识点，意在考查考生灵活运用相关知识综合分析问题的能力。

	（1）设小物块 A 运动到 O 点的速度为 $v_{0}$，由动能定理得

	$mgh = \frac{1}{2}mv_{0}^{2}$

	A、B 发生完全弹性碰撞，设第一次碰撞后瞬间 A 和 B 的速度分别为 $v_{1}$ 和 $v_{2}$，选水平向右为正方向，由动量守恒得

	$mv_{0} = mv_{1} + 4mv_{2}$

	由能量守恒得

	$\frac{1}{2}mv_{0}^{2} = \frac{1}{2}mv_{1}^{2} + \frac{1}{2} \times 4mv_{2}^{2}$

	联立解得 $v_{1} = -\frac{3}{5}\sqrt{2gh}$，$v_{2} = \frac{2}{5}\sqrt{2gh}$，所以第一次碰撞后瞬间 A 和 B 速度的大小分别为 $\frac{3}{5}\sqrt{2gh}$ 和 $\frac{2}{5}\sqrt{2gh}$。

	（2）由题意知，A、B 在地面上的加速度大小相同，均为 $a = \mu g$，

	第一次碰撞后 B 向右运动的位移和时间分别为

	$x_{B} = \frac{v_{2}^{2}}{2a}$，$t_{B} = \frac{v_{2}}{a}$

	解得 $x_{B} = \frac{8}{25}h$，$t_{B} = \frac{4}{5}\sqrt{\frac{2h}{g}}$。

	第一次碰撞后 A 从 O 点向左运动到回到 O 点所用的时间 $t_{A} > \frac{h}{-v_{1}} = \frac{5}{6}\sqrt{\frac{2h}{g}} > t_{B}$，即 A、B 再次碰撞前，B 处于静止状态。

	设 A 向右运动到与 B 相碰前的速度为 $v$，有

	$v^{2} - v_{1}^{2} = -2ax_{B}$

	A、B 发生完全弹性碰撞，设第二次碰撞后瞬间 A 和 B 的速度分别为 $v_{3}$ 和 $v_{4}$，由动量守恒得

	$mv = mv_{3} + 4mv_{4}$

	由能量守恒得

	$\frac{1}{2}mv^{2} = \frac{1}{2}mv_{3}^{2} + \frac{1}{2} \times 4mv_{4}^{2}$

	第二次碰撞后，B 向右运动的位移为 $x_{B}' = \frac{v_{4}^{2}}{2a}$

	第二次碰撞后，A 向左运动的位移为 $x_{A}' = \frac{v_{3}^{2}}{2a}$

	联立解得 $x_{B}' = \frac{8}{125}h$，$x_{A}' = \frac{18}{125}h < x_{B}$。

	即此时 A、B 均静止在水平地面上，A、B 间的距离为

	$x = x_{A}' + x_{B}' = \frac{26}{125}h$。
}

\item
%% number: 15
%% paperName: 2018年海南理综第 15 题 8 分
%% typeId: 6
%% type: 计算题
%% chapter: 气体性质
%% point: 气体连接体
%% method: 无
%% score: 8
%% degree: 450
%% duplicateId: 0
%% body:
（1）如图，一定量的理想气体，由状态 a 等压变化到状态 b，再从 b 等容变化到状态 c，a、c 两状态温度相等。下列说法正确的是\xzanswer{BD}（填入正确答案标号。选对 1 个得 2 分，选对 2 个得 4 分；有选错的得 0 分）

\onepicture[w=4.0cm]{\includesvg{figs/15}}

\fourchoices[answer=BD]
{从状态 b 到状态 c 的过程中气体吸热}
{气体在状态 a 的内能等于在状态 c 的内能}
{气体在状态 b 的温度小于在状态 a 的温度}
{从状态 a 到状态 b 的过程中气体对外做正功}

（2）一储存氮气的容器被一绝热轻活塞分隔成两个气室 A 和 B，活塞可无摩擦地滑动。开始时用销钉固定活塞，A 中气体体积为 $2.5\times10^{-4}\Umc$，温度为 $27\UdC$，压强为 $6.0\times10^{4}\UPa$；B 中气体体积为 $4.0\times10^{-4}\Umc$，温度为 $-17\UdC$，压强为 $2.0\times10^{4}\UPa$。现将 A 中气体的温度降至 $-17\UdC$，然后拔掉销钉，并保持 A、B 中气体温度不变，求稳定后 A 和 B 中气体的压强。

%% answer: 
\jdanswer{
$p = 3.2\times10^{4}\UPa$
}
%% memo:
\memoanswer{
	（1）本题考查对一定质量的理想气体的 $p-V$ 图线的理解、理想气体状态方程、热力学第一定律、理想气体内能及其相关的知识点。

	内能是组成物体分子的无规则热运动动能和分子间相互作用势能的总和，由于理想气体不考虑分子势能，故理想气体的内能等于分子平均动能的总和，而温度是分子平均动能的宏观表现。由理想气体状态方程 $\frac{p_{b}V_{b}}{T_{b}} = \frac{p_{c}V_{c}}{T_{c}}$ 可知，当 $V_{b} = V_{c}$，$p_{b} > p_{c}$ 时，$T_{b} > T_{c}$，故 $\Delta U_{cb} < 0$，根据热力学第一定律 $\Delta U_{cb} = W + Q$，体积 $V$ 不变，故 $W = 0$，所以 $Q < 0$，从状态 b 到状态 c 的过程中气体放热，选项 A 错误；同理，气体在状态 a 的温度等于在状态 c 的温度，故气体在状态 a 的内能等于在状态 c 的内能，选项 B 正确；由理想气体状态方程 $\frac{p_{a}V_{a}}{T_{a}} = \frac{p_{b}V_{b}}{T_{b}}$ 可知，当 $p_{a} = p_{b}$，$V_{a} < V_{b}$ 时，$T_{a} < T_{b}$，选项 C 错误；从状态 a 到状态 b 的过程中气体膨胀对外做正功，故 D 正确。

	（2）本题考查查理定律、玻意耳定律、关联气体、压强及其相关的知识点。

	由题意可知，A 中气体初状态压强为 $p_{A} = 6.0\times10^{4}\UPa$，体积 $V_{A} = 2.5\times10^{-4}\Umc$，温度 $T_{A} = 300\UK$；

	B 中气体初状态压强为 $p_{B} = 2.0\times10^{4}\UPa$，体积 $V_{B} = 4.0\times10^{-4}\Umc$，温度 $T_{B} = 256\UK$。

	拔掉销钉稳定后 A、B 中气体压强相等，设为 $p$。

	设 A 中气体体积为 $V_{A}'$，温度为 $T_{A}' = 256\UK$；设 B 中气体体积为 $V_{B}'$，则

	$V_{A}' + V_{B}' = V_{A} + V_{B} = 6.5\times10^{-4}\Umc$

	对 A 中气体，由理想气体状态方程得

	$\frac{p_{A}V_{A}}{T_{A}} = \frac{pV_{A}'}{T_{A}'}$

	对 B 中气体，由玻意耳定律得

	$p_{B}V_{B} = pV_{B}'$

	联立解得 $p = 3.2\times10^{4}\UPa$。
}

\item
%% number: 16
%% paperName: 2018年海南理综第 16 题 8 分
%% typeId: 6
%% type: 计算题
%% chapter: 光学
%% point: 光的折射
%% method: 无
%% score: 8
%% degree: 550
%% duplicateId: 0
%% body:
（1）警车向路上的车辆发射频率已知的超声波，同时探测反射波的频率。下列说法正确的是\xzanswer{AB}（填入正确答案标号。选对 1 个得 2 分，选对 2 个得 4 分；有选错的得 0 分）

\fourchoices[answer=AB]
{车辆匀速驶向停在路边的警车，警车探测到的反射波频率增高}
{车辆匀速驶离停在路边的警车，警车探测到的反射波频率降低}
{警车匀速驶向停在路边的汽车，探测到的反射波频率降低}
{警车匀速驶离停在路边的汽车，探测到的反射波频率不变}

（2）如图，由透明介质构成的半球壳的内外表面半径分别为 $R$ 和 $\sqrt{2}R$。一横截面半径为 $R$ 的平行光束入射到半球壳内表面，入射方向与半球壳的对称轴平行，所有的入射光线都能从半球壳的外表面射出。已知透明介质的折射率为 $n = \sqrt{2}$。求半球壳外表面上有光线射出区域的圆形边界的半径。不考虑多次反射。

\onepicture[w=5.5cm, align=right]{figs/16.png}

%% answer: 
\jdanswer{
$\frac{\sqrt{3} + 1}{2}R$
}
%% memo:
\memoanswer{
	（1）车辆匀速驶向停在路边的警车，警车探测到的反射波的频率增高，故 A 正确；车辆匀速驶离停在路边的警车，警车探测到的反射波的频率降低，故 B 正确；警车匀速驶向停在路边的汽车，探测到的反射波的频率增高，故 C 错误；警车匀速驶离停在路边的汽车，探测到的反射波的频率降低，故 D 错误。

	（2）本题考查折射定律及其相关的知识点。

	分析边缘光线 a，如下图：

	\onepicture[w=5.0cm]{figs/16_an1.jpg}

	由几何关系得，$\frac{\sin 90^\circ}{\sin\angle CAB} = \sqrt{2}$，可得 $\angle CAB = 45^\circ$，

	在 $\triangle OAB$ 中，$AC = AB$，设为 $r$，在 $\triangle OBC$ 中，由勾股定理有：

	$(R + r_{0})^{2} + r_{0}^{2} = (\sqrt{2}R)^{2}$

	进而求出 $r_{0} = \frac{\sqrt{3} - 1}{2}R$

	故半球壳外表面上有光线射出区域的圆形边界的半径为 $r = r_{0} + R = \frac{\sqrt{3} + 1}{2}R$。

	光路图如图所示，则 $BO'$ 为圆形边界的半径。由折射定律得

	$n = \frac{\sin 90^\circ}{\sin\theta}$，解得 $\theta = 45^\circ$，

	在 $\triangle ABO$ 中，由余弦定理得

	\onepicture[w=4.5cm]{figs/16_an2.jpg}

	$R^{2} + |AB|^{2} - 2R|AB|\cos(180^\circ - \theta) = (\sqrt{2}R)^{2}$，

	解得 $|AB| = \frac{\sqrt{6} - \sqrt{2}}{2}R$，

	$|BO'| = R + |AB|\cos\theta = \frac{\sqrt{3} + 1}{2}R$。
}

\end{enumerate}

\end{document}
